% F -- Where the published ceilings sit, where the one-line rule sits, and how
% thin the band left to work in actually is.
\begin{figure}[!htbp]
\centering
\fitwidth{%
\begin{tikzpicture}[x=0.132cm,y=1cm,font=\small,
  rowlab/.style={anchor=south west,font=\footnotesize,text=clmodel},
  val/.style={anchor=west,font=\footnotesize\bfseries},
]
% --- the scale ----------------------------------------------------------
\draw[clrule] (0,0) -- (100,0);
\foreach \xx in {0,20,40,60,80,100}{
  \draw[clrule] (\xx,0) -- (\xx,-0.14);
  \node[anchor=north,font=\scriptsize,text=clrule] at (\xx,-0.18) {\xx};
}
\node[anchor=north,font=\scriptsize,text=clrule] at (50,-0.58)
  {correct answers out of a hundred};

% --- row 1: the famous ceiling -----------------------------------------
\node[rowlab] at (0,4.26)
  {\textbf{93\,\%} \ the published ceiling: \emph{where will this person
   \textbf{be} in the next hour?}~\cite{song2010limits}};
\fill[clnonsig] (0,3.78) rectangle (93,4.22);
\node[val,text=clrule] at (94,4.00) {93};

% --- row 2: the honest ceiling -----------------------------------------
\node[rowlab] at (0,3.16)
  {\textbf{71\,\%} \ the same ceiling: \emph{where will this person
   \textbf{go} next?}~\cite{ikanovic2017alternative}};
\fill[clnonsig!75] (0,2.68) rectangle (71,3.12);
\draw[clrule,dashed] (71,2.68) rectangle (93,3.12);
\node[anchor=east,font=\footnotesize\bfseries,text=white] at (70,2.90) {71};
\node[font=\scriptsize,text=clrule] at (82,2.90) {22: standing still};

% --- row 3: the rule on a raw day --------------------------------------
\node[rowlab] at (0,2.06)
  {\textbf{69.2\,\%} \ one line of code, on an ordinary day of my records};
\fill[clbase] (0,1.58) rectangle (69.2,2.02);
\node[val,text=clbase] at (70.4,1.80) {69.2};

% --- row 4: the rule and the model, on people who move -----------------
\node[rowlab] at (0,0.96)
  {\textbf{50.2\,\%} \ the same line of code on people who genuinely move, and
   \textbf{54.5\,\%} for my model};
\fill[clbase] (0,0.48) rectangle (50.2,0.92);
\fill[clgain] (50.2,0.48) rectangle (54.5,0.92);
\node[val,text=clgain] at (55.7,0.70) {54.5};

% --- the reading --------------------------------------------------------
\node[draw=clrule,fill=clsoft,rounded corners=3pt,align=left,inner sep=7pt,
      text width=13.0cm,font=\footnotesize,anchor=north west] at (0,-0.98)
  {\textbf{Why the third row is the problem.} The one line of code answers
   \emph{``the same antenna as a moment ago''}, every time. At 69.2 against an
   honest ceiling of 71 there is nothing at all left to win: on a day file taken
   as it comes, one line of code is already at the ceiling. That is why every
   result in this report is measured on people who genuinely move, where the
   same line of code drops to 50 and a model has room to show.};
\end{tikzpicture}}
\caption{The room a predictor has, once standing still is taken out of it}
\label{fig:ceiling-ladder}
\figtext{%
\figpara{Reading the four bars of Figure~\ref{fig:ceiling-ladder}.}{All four bars in Figure~\ref{fig:ceiling-ladder} are on one scale, the share of moments
answered correctly. The two grey bars are published ceilings, the best score any
predictor could ever reach; the two blue ones are what the trivial rule and fine-tuned
model score on the evaluation data. Song et al.\ \textbf{\cite{song2010limits}} put the ceiling at 93\,\%. Ikanovic and
Mollgaard \textbf{\cite{ikanovic2017alternative}} re-ran the same measurement after deleting the moments where the person
had not moved, and it fell to 71\,\%, the dashed part of the second bar in Figure~\ref{fig:ceiling-ladder}.}

\figpara{Comparison of bottom bars across dataset cohorts in Figure~\ref{fig:ceiling-ladder}.}{An
ordinary day file and a selection of people who genuinely move are not comparable
on absolute score in Figure~\ref{fig:ceiling-ladder}. The rule scores 69.2 on the first and 50.2 on the second
without changing a character of its own definition; \textbf{Chapter~\ref{ch:data}} makes
that move and explains it.}}

\end{figure}
